% © 2025 Ing. David Jaroš — CC BY-NC-ND 4.0
%
% This work is licensed under a Creative Commons Attribution-NonCommercial-NoDerivatives
% 4.0 International License (CC BY-NC-ND 4.0).

\documentclass[11pt,a4paper]{article}
\usepackage{amsmath, amssymb, amsthm}
\usepackage{hyperref}
\usepackage{geometry}
\geometry{margin=2.5cm}

\newtheorem{theorem}{Theorem}
\newtheorem{definition}{Definition}
\newtheorem{proposition}{Proposition}
\newtheorem{remark}{Remark}

\title{Canonical Metric Definition\\
\large Unified Biquaternion Theory — Canonical Theory Series}
\author{Ing.\ David Jaroš}
\date{2025}

\begin{document}
\maketitle

\begin{abstract}
This document gives the canonical definition of the metric tensor in UBT.
The key point is that the classical spacetime metric $g_{\mu\nu}$ is
\emph{derived} from the biquaternionic field $\Theta$, not postulated.
Two equivalent derivations are given and their equivalence is proved.
\end{abstract}

\section{The Biquaternionic Metric}

\begin{definition}[Biquaternionic Metric Tensor]\label{def:biqmetric}
The biquaternionic metric tensor is:
\begin{equation}\label{eq:biqmetric}
  \mathcal{G}_{\mu\nu}[\Theta]
  \;\in\; \mathbb{H}\otimes_\mathbb{R}\mathbb{C},
\end{equation}
constructed from $\Theta$ as the quadratic form:
\begin{equation}\label{eq:Gfromtheta}
  \mathcal{G}_{\mu\nu}
  \;=\;
  \mathrm{Tr}\bigl(\partial_\mu\Theta\cdot(\partial_\nu\Theta)^\dagger\bigr).
\end{equation}
It encodes the full biquaternionic geometry including real and imaginary parts.
\end{definition}

\noindent\textbf{Source}: \texttt{canonical/geometry/biquaternion\_metric.tex}, \S1--\S2.

\bigskip
The biquaternionic metric decomposes as:
\begin{equation}
  \mathcal{G}_{\mu\nu}
  \;=\;
  \underbrace{g_{\mu\nu}}_{\text{real (GR)}}
  \;+\;
  i\,\underbrace{h_{\mu\nu}}_{\text{phase curvature}}
  \;+\;
  \underbrace{\mathbf{k}_{\mu\nu}}_{\text{quaternionic}},
\end{equation}
where $\mathbf{k}_{\mu\nu}$ contains the quaternionic (non-abelian) components.

\section{Classical Metric: Real Projection}

\begin{definition}[Physical Spacetime Metric]\label{def:metric}
The classical spacetime metric is the real part of the biquaternionic metric:
\begin{equation}\label{eq:gRe}
  g_{\mu\nu}(x)
  \;:=\;
  \mathrm{Re}\bigl(\mathcal{G}_{\mu\nu}[\Theta]\bigr)
  \;=\;
  \mathrm{Re}\,\mathrm{Tr}\bigl(\partial_\mu\Theta\cdot(\partial_\nu\Theta)^\dagger\bigr).
\end{equation}
\textbf{Forbidden}: $g_{\mu\nu}$ must not be introduced without explicit reference
to the parent $\mathcal{G}_{\mu\nu}$.
\end{definition}

\noindent\textbf{Source}: \texttt{canonical/geometry/metric.tex}, \S2; Rule M1.

\section{Equivalent Tetrad-Based Definition}

\begin{definition}[Tetrad Form]\label{def:tetrad}
Equivalently, introducing tetrads $E_\mu := \partial_\mu\Theta$:
\begin{equation}\label{eq:gtetrad}
  g_{\mu\nu}
  \;=\;
  \mathrm{Re}\,\mathrm{Sc}\bigl(E_\mu\,E_\nu^\dagger\bigr),
\end{equation}
where $\mathrm{Sc}(\cdot)$ denotes the scalar (real quaternion) part.
\end{definition}

\begin{theorem}[Equivalence of Metric Definitions]
Definitions~\ref{def:metric} and \ref{def:tetrad} are equivalent:
\begin{equation}
  \mathrm{Re}\,\mathrm{Tr}\bigl(\partial_\mu\Theta\cdot(\partial_\nu\Theta)^\dagger\bigr)
  \;=\;
  \mathrm{Re}\,\mathrm{Sc}\bigl(\partial_\mu\Theta\cdot(\partial_\nu\Theta)^\dagger\bigr),
\end{equation}
with $E_\mu = \partial_\mu\Theta$.
\end{theorem}

\begin{proof}
Under the identification $E_\mu = \partial_\mu\Theta$, both definitions yield
$\mathrm{Re}(E_\mu E_\nu^\dagger)$; the trace and scalar-part operations coincide
on the real component. Proved in detail in
\texttt{consolidation\_project/GR\_closure/step1\_metric\_bridge.tex}.
\end{proof}

\noindent\textbf{Status}: PROVED [L0].

\section{Signature and Non-Degeneracy}

\begin{proposition}[Lorentzian Signature]
The projected metric $g_{\mu\nu}$ has Lorentzian signature $(-,+,+,+)$
(or equivalently $(+,-,-,-)$ in the opposite convention) for physically
admissible configurations of $\Theta$.
\end{proposition}

\noindent\textbf{Status}: PROVED [L0].\\
\noindent\textbf{Source}: \texttt{consolidation\_project/GR\_closure/step3\_signature\_theorem.tex}.

\begin{proposition}[Non-Degeneracy]
For all non-trivial field configurations $\Theta \not\equiv 0$,
the projected metric $g_{\mu\nu}$ is non-degenerate: $\det(g_{\mu\nu}) \neq 0$.
\end{proposition}

\noindent\textbf{Status}: PROVED [L0].\\
\noindent\textbf{Source}: \texttt{consolidation\_project/GR\_closure/step2\_nondegeneracy.tex}.

\section{Phase Projection Generalization}

For a general constant phase $\varphi$, the projected metric is:
\begin{equation}
  g_{\mu\nu}^{(\varphi)}
  \;=\;
  \mathrm{Re}\bigl(e^{-i\varphi}\,\mathcal{G}_{\mu\nu}\bigr).
\end{equation}
The choice $\varphi = 0$ is a gauge choice (phase frame), not a physical
restriction. Different phase frames are related by U(1) transformations
(phase frame bundle).

\noindent\textbf{Source}: \texttt{canonical/geometry/phase\_projection.tex}, \S2.

\section{Connection to GR}

The metric defined by eq.\ \eqref{eq:gRe} satisfies:
\begin{enumerate}
  \item Symmetry: $g_{\mu\nu} = g_{\nu\mu}$ (follows from definition).
  \item Lorentzian signature: $\mathrm{sign}(g) = (-,+,+,+)$ (Proved [L0]).
  \item Non-degeneracy: $\det(g)\neq 0$ (Proved [L0]).
  \item Einstein equations: $G_{\mu\nu} = \kappa T_{\mu\nu}$ upon variation of
        the action (Proved [L1]; see \texttt{canonical/bridges/GR\_chain\_bridge.tex}).
\end{enumerate}

\section{Cross-References}

\begin{itemize}
  \item Full biquaternionic metric: \texttt{canonical/geometry/biquaternion\_metric.tex}
  \item Phase projection: \texttt{canonical/geometry/phase\_projection.tex}
  \item Metric equivalence proof: \texttt{consolidation\_project/GR\_closure/step1\_metric\_bridge.tex}
  \item Non-degeneracy: \texttt{consolidation\_project/GR\_closure/step2\_nondegeneracy.tex}
  \item Signature: \texttt{consolidation\_project/GR\_closure/step3\_signature\_theorem.tex}
  \item GR chain: \texttt{canonical/bridges/GR\_chain\_bridge.tex}
\end{itemize}

\end{document}
